本节按论文章节顺序列出全部符号。“含义”一栏括号内是该符号在原文中首次出现的行号；“实证中的处理”一栏说明论文第 5 节如何对待该量，分为直接观测或估计、用代理量、只给情景区间、不可识别四类。表~\ref{tab:sym-model}、表~\ref{tab:sym-design} 与表~\ref{tab:sym-emp} 依次对应模型设定、估计与设计、预测检验与实证三部分；注明“本讲义引入”的符号是为补全推导而设，原文没有。

量纲约定如下。理论部分的一期是一个传播代，实证部分的一期是一周（欧洲 COVID-19 数据为一天）。$R_t$ 在理论中是每名感染者的期望继发感染数，在实证中由周增长比 $G$ 代理，两者只有在代际间隔约为一周时才可直接比较，第~\ref{sec:checks} 节讨论这一口径差的后果。

\begin{table}[!tbp]
\centering
\footnotesize
\caption{传播与报告模型的符号（原文第 2 节）。}
\label{tab:sym-model}
\begin{tabular}{L{1.7cm}L{5.8cm}L{3.0cm}L{4.4cm}}
\toprule
符号 & 含义（首次出现行） & 取值与单位 & 实证中的处理 \\
\midrule
$t$ & 传播代（期）下标（89） & $t=1,2,\dots$ & 对应州周 \\
$i$ & 上一代第 $i$ 名感染者（89） & $1,\dots,n$ & 无 \\
$n$ & 上一传播代的真实感染者数，即真实风险集规模（89） & 正整数；人 & 不可观测，用 $D_t\approx\rho n$ 间接代理 \\
$X_{i,t}$ & 第 $i$ 名感染者产生的继发感染数（89） & 非负整数；人 & 不可观测 \\
$R_t$ & 给定当期环境的传播参数，即个体期望继发感染数，目标 A（89） & $R_t>0$；无量纲 & 无直接估计，以 $G_{j,t+1}=C_{j,t+1}/D_{j,t}$ 作经验代理 \\
$k$ & 负二项离散参数，刻画个体传播异质性（92） & $k>0$；$k\to0$ 为极端超传播，$k\to\infty$ 退化为 Poisson & 没有估计；情景区间 $[0.2,0.5]$（515 行），仿真取 0.1、0.35、1.0 \\
$I_t$ & 当期真实新发感染总数 $\sum_iX_{i,t}$（94） & 非负整数；人 & 不可观测 \\
$\rho$ & 报告率，即每个真实感染被报告的概率（96） & $(0,1]$ & 没有估计；情景区间 $[0.0056,0.021]$（515 行），医院填报覆盖率作制度层面的代理（467 行） \\
$C_t$ & 当期报告病例数（98） & 非负整数；例 & 直接观测（NHSN 流感住院数） \\
$\bar R$ & 跨期基础传播水平 $\mathbb{E}[R_t\mid n]$，目标 B（107） & $>0$ & 阶段内 $G$ 的样本平均：初期 1.6706，达峰期 0.8075 \\
$v_R$ & 不可约环境方差 $\operatorname{Var}(R_t\mid n)$（111） & $>0$；无量纲平方 & 不可单独识别（截距 $a$ 混入其他与规模无关的误差）；情景区间 $[0.04,0.08]$，仿真取 0.04 \\
$\mathcal F_t^I$ & 全历史信息集 $\sigma(I_{1:t},C_{1:t})$（120） & $\sigma$-代数 & 无 \\
$\mathcal F_t^C$ & 仅含报告的信息集 $\sigma(C_{1:t})$（121） & $\sigma$-代数，$\subset\mathcal F_t^I$ & 实际可用的信息 \\
$\Lambda_{t+1}$ & 下一期感染压力 $\sum_sw_sI_{t+1-s}$（120，定义在 291） & $\ge0$；人 & 不可观测，由 $D_t$ 代理 \\
$w_s$ & 代际滞后权重（291） & $w_s\ge0$，$\sum_sw_s=1$ & 实证取 $(0.65,0.25,0.10)$，敏感性取 $(1/3,1/3,1/3)$ 与 $(1,0,0)$；仿真取 $(0.45,0.35,0.15,0.05)$ \\
$S$ & 滞后窗口长度（291） & 正整数；期 & 实证 $S=3$，仿真 $S=4$ \\
$\lambda_i$ & 个体传染力，$\lambda_i\mid R_t\sim\mathrm{Gamma}(k,R_t/k)$（本讲义引入） & $>0$ & 无 \\
$p$ & 个体负二项的成功概率 $k/(k+R_t)$（688） & $(0,1)$ & 无 \\
$p^*$ & 抽稀后负二项的成功概率 $k/(k+\rho R_t)$（693） & $(0,1)$ & 无 \\
$G_{I_t}(s)$, $G_{C_t}(s)$ & 概率生成函数（688） & $\lvert s\rvert\le1$ & 无 \\
\bottomrule
\end{tabular}
\end{table}

\begin{table}[!tbp]
\centering
\footnotesize
\caption{估计量、方差、交叉点与预算设计的符号（原文第 3.1 至 3.5 节）。}
\label{tab:sym-design}
\begin{tabular}{L{1.7cm}L{5.8cm}L{3.0cm}L{4.4cm}}
\toprule
符号 & 含义（首次出现行） & 取值与单位 & 实证中的处理 \\
\midrule
$\widehat R_t$ & 单期估计量 $C_t/(\rho n)$，同时是最大似然估计（135） & $\ge0$ & 实证对应物是 $G=C_{t+1}/D_t$ \\
$\mathcal I(R_t\mid n)$ & 关于 $R_t$ 的 Fisher 信息（138） & $>0$；量纲为 $1/R_t^2$ & 无 \\
$m$ & 期望报告规模 $\rho n$（147，定义在 174） & $>0$；例/期 & 用 $D_t$ 代理 \\
$A$ & 规模相关方差系数 $\bar R+\rho(\bar R^2+v_R)/k$（175） & $>\bar R$；量纲为例，使 $A/m$ 无量纲 & 无分母噪声时理论上对应 $b\approx A/\bar R^2$ \\
$m_\times$ & 方差占比交叉点 $A/v_R$（164，定义在 183） & $>0$；例/期 & 经验值 $b/a$，初期 22.6 [13.7, 33.7] 例/周 \\
$V(n,\rho)$ & 把原文式 (7) 写成 $(n,\rho)$ 的二元函数（204） & $>v_R$ & 无 \\
$T$ & 观测期数（237） & 正整数；推论 1 中连续放松为 $T\ge1$ & 实证未使用 \\
$n_T$ & 每期相同的真实风险集规模 $B/(\rho T)$（237） & 正整数 & 无 \\
$\mathcal D_T$ & 设计方案 $(T,n_T)$（237） & 无 & 无 \\
$\overline{\widehat R}_T$ & 跨期样本平均估计量 $\frac1T\sum_t\widehat R_t$（241） & 无 & 无 \\
$B$ & 总规模预算 $Tm$，即各期期望被报告人数之和（245） & 例 & 无 \\
$r$ & 环境扰动的 AR(1) 自相关系数（252） & $0\le r<1$ & 没有估计；可由 $G$ 的周自相关估计 \\
$T_{\rm eff}$ & 有效独立期数 $T(1-r)/(1+r)$（254） & $\le T$ & 无 \\
$\mathcal T(B,\rho)$ & 整数可行期数集合（259） & $\subset\mathbb N_+$ & 无 \\
$c_0$ & 每期固定运转成本（274） & $>0$；货币/期 & 没有估计，只给成本比情景 $c_0/c_1\in\{10,100\}$ \\
$c_1$ & 每报告一例的边际成本（274） & $>0$；货币/例 & 同上 \\
$c_1'$ & 按覆盖人群计的人均成本，$c_1=c_1'/\rho$（274） & 货币/人 & 无 \\
$\mathcal B$ & 总经费预算 $T(c_0+c_1m)$（274） & 货币 & 无 \\
$m^*$ & 最优单期规模 $\sqrt{(c_0/c_1)m_\times}$（277） & 例/期 & 情景值 14.6--22.2（$c_0/c_1=10$）、46.1--70.2（$c_0/c_1=100$） \\
$v_{\rm eff}$ & AR(1) 下的有效环境方差 $v_R(1+r)/(1-r)$（280） & 无 & 无 \\
$m^*_{\rm AR}$ & AR(1) 修正的最优规模（280） & 例/期 & 无 \\
\bottomrule
\end{tabular}
\end{table}

\begin{table}[!tbp]
\centering
\footnotesize
\caption{预测、检验与实证的符号（原文第 3.6 节至第 5 节）。}
\label{tab:sym-emp}
\begin{tabular}{L{1.9cm}L{5.8cm}L{2.8cm}L{4.4cm}}
\toprule
符号 & 含义（首次出现行） & 取值与单位 & 实证中的处理 \\
\midrule
$H_0$, $H_1$ & 原假设 $R_t\le1$ 与备择 $R_t>1$（318） & 无 & 实证未做 \\
$\alpha$ & 名义检验水平（324） & 0.05 & 无 \\
$c_\alpha$ & 离散临界值，即使 $\Pr_{R=1}(C\ge c)\le\alpha$ 的最小整数（326） & 非负整数 & 无 \\
$\delta$ & 备择增长强度，$R_t=1+\delta$（328） & $>0$ & 仿真取 0.15、0.25、0.30、0.50 \\
$\pi(n,\delta)$ & 精确检出功效（330） & $[0,1]$ & 无 \\
$F_{\rm NB}(x;r,p)$ & 负二项累积分布函数，此处 $r$ 是形状参数（332） & 无 & 与 AR(1) 的 $r$ 撞名 \\
$Z$, $V_0$, $z_{1-\alpha}$ & Wald 统计量、原假设下方差 $1/(\rho n)+1/(kn)$、正态分位数（350） & 无 & 无 \\
$D_t$ & 历史报告加权和 $\sum_sw_sC_{t+1-s}$（373） & $\ge0$；例 & 由数据直接算出，$D_t\ge5$ 入样 \\
$G_{t+1}$ & 事后增长统计量 $C_{t+1}/D_t$（379） & $\ge0$ & 直接计算 \\
$\widehat C_{t+1\mid t}$ & 事前单步预测 $\bar RD_t$（380） & 例 & 实证为 $\bar R_{\rm phase}D_{j,t}$ \\
$\operatorname{RelMSE}(m)$ & 相对均方预测误差（382） & $\ge0$ & 分箱计算 \\
$\mu_{t+1\mid t}$, $\sigma^2_{t+1\mid t}$ & $C_{t+1}$ 在 $\mathcal F_t^C$ 下的条件期望与条件方差（387） & 无 & 无 \\
$j$, $J$ & 辖区下标与辖区数（443） & 50 州加哥伦比亚特区，$J=51$ & 无 \\
$D_{j,t}$, $G_{j,t+1}$ & 州 $j$ 第 $t$ 周的历史有效输入压力与增长代理（443、445） & 例/周；无量纲 & 直接计算 \\
$\bar R_{\rm early}$, $\bar R_{\rm peak}$ & 初期与达峰期的 $G$ 样本平均（455、558） & 1.6706、0.8075 & 回顾性样本平均；滚动口径见原文第 551 行 \\
$a$ & 仿射模型截距，即与规模无关的经验损失下限（506） & 初期 0.0991 [0.0841, 0.1173] & FGLS 估计 \\
$b$ & 仿射模型斜率，即规模相关项（506） & 初期 2.2434 [1.5019, 3.0161] & FGLS 估计 \\
$L$ & 单观测损失 $((G-\bar R)/\bar R)^2$（512） & $\ge0$ & 直接计算 \\
$\hat\mu$ & FGLS 第一步拟合值，权重取 $1/\hat\mu^2$（512） & $>0$ & 无 \\
$D_t^{\rm iv}$ & 工具变量，取滞后 3--5 周的规模（515） & 例/周 & 无 \\
$b_{\rm iv}$, $b_{\rm ols}$ & 两阶段最小二乘与 OLS 斜率（515） & 无 & 无 \\
$\tau$ & Cori 方法的滑动窗长度（519） & 1、2、3 周 & 无 \\
$\bar c$ & 辖区间环境扰动的平均空间相关系数（617） & $[0,1)$ & 没有估计 \\
$v_R^{\rm agg}$ & 汇总后的环境方差 $v_R[\bar c+(1-\bar c)/J]$（617） & 无 & 无 \\
$\widehat Y_{j,t}$ & FluSight 集成模型发布的预期住院规模（624） & 例/周 & 替代 $D_t$ 作自变量 \\
\bottomrule
\end{tabular}
\end{table}

\subsection{撞名与一词多义}

论文中有几处同一符号承担两种含义，推导时容易混淆，下面逐条说明。

\begin{itemize}
\item $A$ 同时用于目标 A（估计 $R_t$）和方差系数 $A$。原文第 251 行在命题 4 中作了区分，其余位置没有说明。
\item $r$ 同时用于 AR(1) 自相关系数和 $F_{\rm NB}(x;r,p)$ 的形状参数，原文第 418 行还写了 $r=nk$。
\item $B$ 是总规模预算，单位为例；$\mathcal B$ 是总经费，单位为货币。
\item $m$ 在理论中是非随机的期望值 $\rho n$，在实证中由随机的已实现量 $D_t$ 代理，论文图 3 的横轴又是初始期望报告规模（原文第 375、393 行）。
\item $k$ 是个体离散参数，$C_t$ 的负二项形状参数是 $nk$。
\item $p$ 与 $p^*$ 分别是抽稀前后的成功概率。
\item 正文的第 5 个命题是预测分解、第 6 个是增长检验，附录小节标题的顺序与之相反。本讲义一律采用正文编号。
\end{itemize}
